\documentclass[11pt,a4paper]{article}

\usepackage[utf8]{inputenc}
\usepackage[T1]{fontenc}
\usepackage[spanish]{babel}
\usepackage[margin=2.5cm]{geometry}

\title{Pipeline automatizado de IaC con Terraform \\ \large Despliegue de un Data Lake en Google Cloud Platform}
\author{
  Ana Inado \\
  Ariana Sanchez \\
  Wilber Correa \\
  Linder Kriss \\
  Brenner Perez
}
\date{14 de julio del 2026}

\begin{document}

\maketitle

\section{Introducción y caso de uso}

Este trabajo diseña e implementa el esquema de un pipeline CI/CD automatizado, basado en Infraestructura como Código (IaC) con Terraform, para desplegar un servicio de datos en la nube. Se seleccionó \textbf{Google Cloud Platform (GCP)} como proveedor y un \textbf{Data Lake} como caso de uso: tres buckets de Cloud Storage (\texttt{raw}, \texttt{processed}, \texttt{curated}), un dataset de BigQuery y una cuenta de servicio con permisos para Dataflow, declarados en \texttt{terraform/*.tf} y orquestados mediante \texttt{cloudbuild.yaml}.

Código fuente completo: \texttt{https://github.com/wilberever/MaestriaDevops}. El esquema visual está en \texttt{pipeline\_cicd\_gcp.drawio}.

\section{Funcionamiento del pipeline}

El pipeline tiene seis etapas:

\begin{enumerate}
  \item \textbf{Control de versiones.} Repositorio Git con un archivo \texttt{CODEOWNERS} que exige revisión antes de mezclar cambios a \texttt{main}.
  \item \textbf{Integración continua + DevSecOps.} Cada \emph{push} valida sintaxis (\texttt{terraform fmt}, \texttt{terraform validate}) y corre \texttt{tfsec} sobre el código de infraestructura antes de generar el plan. La seguridad es una condición para avanzar, no una revisión final.
  \item \textbf{Planificación.} \texttt{terraform plan} calcula los cambios contra un \emph{backend} remoto en Cloud Storage con bloqueo de estado.
  \item \textbf{Aprovisionamiento.} \texttt{terraform apply} crea los buckets, el dataset de BigQuery, la cuenta de servicio de Dataflow y la política de alertas.
  \item \textbf{Validación.} Un \emph{smoke test} confirma que los tres buckets y el dataset existen y son accesibles.
  \item \textbf{Monitoreo.} Cloud Monitoring y Cloud Alerting notifican por correo ante errores.
\end{enumerate}

\section{Decisiones y una inconsistencia detectada}

Se usa \emph{backend} remoto con versionado (en vez de estado local) porque el archivo \texttt{tfstate} concentra el detalle completo de la infraestructura y, sin protegerlo, es un punto único de fuga de información (Red Hat, s.f.).

El esquema visual del pipeline incluye una aprobación manual antes de \texttt{apply}, pero el código actual ejecuta \texttt{terraform apply -auto-approve} apenas termina el \texttt{plan}, sin ninguna pausa humana. Se documenta esta diferencia en vez de ocultarla: el pipeline hoy prioriza velocidad y confía solo en la validación y en \texttt{tfsec} como filtro antes de tocar infraestructura real, lo cual es razonable en \emph{staging} pero insuficiente para producción.

Tampoco existe todavía, dentro de este mismo repositorio, un job real de Dataflow (Apache Beam): lo que hay es la cuenta de servicio y los permisos para correrlo, no el \emph{pipeline} ETL en sí.

\section{Conclusiones y recomendaciones}

El pipeline cubre las seis etapas exigidas, integrando la seguridad dentro del flujo (DevSecOps) en vez de como revisión final, en línea con lo señalado por IBM (s.f.) y Red Hat (s.f.). La automatización reduce el error humano y mejora la trazabilidad, pero también puede propagar un error de configuración más rápido que un proceso manual si faltan controles antes de \texttt{apply}.

Recomendaciones concretas para las próximas iteraciones:

\begin{enumerate}
  \item Agregar una aprobación manual real antes de \texttt{apply} en producción.
  \item Cifrar y restringir el acceso al \emph{bucket} del \texttt{tfstate}.
  \item Prohibir cambios manuales en la consola de GCP y detectar \emph{drift} con un \texttt{terraform plan} programado.
  \item Completar el job real de Dataflow antes de confiar en la alerta de monitoreo que hoy lo supone.
\end{enumerate}

\section*{Referencias}

IBM. (s.f.). \emph{¿Qué es DevOps?} IBM Think. \texttt{https://www.ibm.com/mx-es/think/topics/devops}

Microsoft Azure. (s.f.). \emph{¿Qué es DevOps?} \texttt{https://azure.microsoft.com/es-es/resources/cloud-computing-dictionary/what-is-devops}

Red Hat. (s.f.). \emph{CI/CD: La integración y la distribución continuas.} \texttt{https://www.redhat.com/es/topics/devops/what-is-ci-cd}

HashiCorp. (s.f.). \emph{Terraform documentation.} \texttt{https://developer.hashicorp.com/terraform/docs}

Google Cloud. (s.f.). \emph{Cloud Build documentation.} \texttt{https://cloud.google.com/build/docs}

\end{document}
